\scanpage{124}
\mspage{- 120 -}\hfill IX

un accouplement
\[
 v_\ell^{K'}:\underline M_K\otimes\underline M'_K\longrightarrow\mathbb Z_{\ell K}\quad.\tag*{(*)}
\]

Ce dernier à priori dépend du choix de $K'$, mais il résulte de (10.4.3) ci-dessous que l'accouplement
\[
 u_\ell:\underline M_K\otimes\underline M'_K\longrightarrow\mathbb Q_{\ell K}\tag*{(9.7.1)}
\]
donné par
\[
 u_\ell=(1/e(S'/S))\,v_\ell^{K'}
\]
n'en dépend pas, où $e(S'/S)=\operatorname{long}(V'/\underline m V')$. On pourra appeler encore (9.7.1) l'\uline{accouplement de monodromie} pour le schéma abélien $A_K$ sur $K$ et le nombre premier $\ell$ envisagé, puisque dans le cas où $A_K$ lui-même a réduction semi-stable, il coïncide évidemment avec la restriction à $\eta$ de l'accouplement de monodromie (9.1.2). Il résulte de 10.4 plus bas que cet accouplement est en fait toujours à valeurs dans $\mathbb Q_K$, et est indépendant de $\ell$.

On peut se demander s'il est à valeurs dans $\mathbb Z_K$, ce qui revient à dire que (9.7.1) est à valeurs dans $\mathbb Z_{\ell K}$ pour tout nombre premier $\ell$. C'est le cas du moins si $A_K$ est une courbe elliptique, le seul cas qui demande une vérification étant alors celui où $\underline M_K\simeq\underline M'_K$ est de rang 1, de sorte que $\underline M_K\otimes\underline M'_K$ est canoniquement isomorphe à $\mathbb Z_K$, et que l'accouplement de monodromie (9.7.1) peut s'interprèter simplement comme un nombre rationnel ordinaire. On trouve grâce aux résultats de NERON [19] que celui-ci est égal à $v(j)$, $j$ étant l'invariant absolu de $A_K$ et $v$ la valuation normalisée de $K$.

\mspage{\handwritten{438}}
% stamp: numbered label 438
% archive sheet number: 123
